\documentclass[12pt,a4paper]{ctexart}

\usepackage[a4paper,margin=2.5cm]{geometry}
\usepackage{amsmath,amssymb,amsthm}
\usepackage{array}
\usepackage{booktabs}
\usepackage{enumitem}
\usepackage{float}
\usepackage{graphicx}
\usepackage{hyperref}
\usepackage{longtable}
\usepackage{multirow}
\usepackage{xcolor}

\hypersetup{
  colorlinks=true,
  linkcolor=blue,
  citecolor=blue,
  urlcolor=blue
}

\title{数理逻辑课程大作业实验报告}
\author{组长姓名 \and 组员姓名1 \and 组员姓名2 \and 组员姓名3 \and 组员姓名4}
\date{\today}

\begin{document}

\maketitle

\begin{abstract}
本文报告围绕 NesyLink 数理逻辑课程项目展开，目标是在 Zelda-like 地牢环境中设计可运行的 Agent，并对环境或策略中的关键部分进行 Lean 形式化建模与证明。

本项目已对 \texttt{task\_1} 至 \texttt{task\_5} 全部完成 Lean 环境形式化与可完成性证明，
对应文件位于 \texttt{student\_agent/lean/} 目录下。五个形式化模型均通过 Lean LSP 诊断（零错误），
并遵循统一证明骨架：局部前置/后置引理 $\rightarrow$ 必要性引理 $\rightarrow$
\texttt{Reachable} 归纳不变式 $\rightarrow$ \texttt{completion\_postcondition} $\rightarrow$
witness 计划 $\rightarrow$ 可完成性。其中 \texttt{task\_4} 因涉及旋转桥、条件门、
装备获取与事件驱动显现机制，作为重点展开对象。
\end{abstract}

\tableofcontents
\newpage

\section{项目背景与任务要求}

\subsection{课程作业目标}

本次课程作业要求小组围绕给定游戏环境完成三部分内容：
\begin{enumerate}[leftmargin=2em]
  \item 设计并实现一个能够完成任务的 Agent；
  \item 使用 Lean 对环境或策略中的关键部分进行形式化建模；
  \item 证明若干性质，例如动作合法性、安全性、可达性或策略正确性。
\end{enumerate}

\subsection{环境限制与评测要求}

\begin{itemize}[leftmargin=2em]
  \item 训练和调试阶段可以使用像素、reward 和 \texttt{info} 等信息。
  \item 最终测试阶段，Agent 不能直接依赖环境内部 \texttt{info}。
  \item 最终测试阶段，Agent 只能根据像素观测、reward 和物品栏信息进行决策。
  \item 如果使用机器学习方法，需要说明训练方式、信息来源以及与最终评测接口的一致性。
\end{itemize}

\subsection{本文结构}

本文后续将依次介绍环境与任务分析、总体方法设计、视觉与状态抽取、策略与规划、
Lean 形式化与证明（覆盖 \texttt{task\_1} 至 \texttt{task\_5}）、
实验结果以及项目分工。

\section{环境与任务分析}

\subsection{NesyLink 环境概述}

NesyLink 是一个 Zelda-like 的离散动作地牢环境。默认任务接口采用 Gym 风格：
\texttt{obs, info = env.reset(seed)}，以及
\texttt{obs, reward, terminated, truncated, info = env.step(action)}。
本次课程任务要求最终提交版策略在推理阶段只能使用像素观测、reward
和显式提供的物品栏信息，不能直接读取 \texttt{info} 中的隐藏状态。

动作空间为 7 个离散动作：
\texttt{WAIT}、\texttt{UP}、\texttt{DOWN}、\texttt{LEFT}、\texttt{RIGHT}、
\texttt{BUTTON\_A}、\texttt{BUTTON\_B}。默认像素控制下，房间地图区域大小为
\texttt{160 \(\times\) 128} 像素，对应 \texttt{10 \(\times\) 8} 个 tile，
每个 tile 为 \texttt{16 \(\times\) 16} 像素。默认观测 \texttt{obs}
为 RGB 图像，shape 为 \texttt{(128, 160, 3)}。环境还提供结构化
\texttt{info} 作为调试与监督接口，其中包含房间编号、角色坐标、事件记录、
动态对象状态和奖励信号等信息。

五个任务的差异主要体现在关键机制上：\texttt{task\_1}
是静态单房间取钥匙开门；\texttt{task\_2} 引入怪物、战斗和条件门；
\texttt{task\_3} 要求跨房间记忆；\texttt{task\_4}
引入旋转桥、条件门、装备获取和隐藏宝箱，是最适合做显式环境建模的一关；
\texttt{task\_5} 则是四房间按钮门与全宝箱收集的综合任务。基于课程要求与当前项目进度，
本文对 \texttt{task\_1} 至 \texttt{task\_5} 全部进行了 Lean 环境形式化与可完成性证明，
其中 \texttt{task\_4} 因机制最复杂作为重点展开对象。

\subsection{任务拆解}

建议按任务逐一描述：
\begin{itemize}[leftmargin=2em]
  \item Task 1：单房间取钥匙开门。
  \item Task 2：怪物处理与取钥匙。
  \item Task 3：多房间状态记忆与往返。
  \item Task 4：桥与按钮机制、复杂子任务链。
  \item Task 5：综合关卡与泛化挑战。
\end{itemize}

\subsection{主要难点}

\begin{enumerate}[leftmargin=2em]
  \item 最终测试不能直接读取 \texttt{info}。
  \item 需要从像素中抽取对策略有用的状态表示。
  \item 需要在短时间内兼顾实现、证明和报告。
\end{enumerate}

\section{总体方法设计}

\subsection{系统框架}

建议在此给出整体框图，说明以下模块之间的关系：
\begin{itemize}[leftmargin=2em]
  \item 视觉/状态抽取模块
  \item 策略与规划模块
  \item 安全约束或合法动作筛选模块
  \item 形式化验证模块
\end{itemize}

\subsection{为什么选择该路线}

此处说明为什么选择当前路线，例如：
\begin{itemize}[leftmargin=2em]
  \item 可解释性更强；
  \item 更容易与 Lean 证明对齐；
  \item 更适合一周内完成；
  \item 更容易兼顾性能与可验证性。
\end{itemize}

当前阶段已经选择“统一 Policy 骨架 + 分任务逐步扩展 + 可验证层优先”的路线。
具体而言，项目先在自有代码目录中建立统一的 Agent 入口，
优先跑通 \texttt{task\_1} 的 baseline，再逐步扩展到 \texttt{task\_2} 至 \texttt{task\_5}，
并以阶段机、合法动作筛选器和动态对象状态跟踪作为后续形式化重点。
目前自有 baseline 已覆盖 \texttt{task\_1}、\texttt{task\_2}、\texttt{task\_3} 和 \texttt{task\_4}。
其中 \texttt{task\_2} 仍属于调试用途的脚本化版本，
而 \texttt{task\_3} 已经采用房间级规则状态机，
\texttt{task\_4} 已经采用显式阶段机并维护桥状态与关键事件记忆，
后续需要将这些基线逐步替换为更通用、更加符合最终测试限制的策略。

\section{视觉与状态抽取}

\subsection{输入与输出设计}

此处描述从像素观测中提取哪些信息，例如玩家位置、墙体、钥匙、怪物、按钮、门、桥状态等。

\subsection{实现方法}

此处记录：
\begin{itemize}[leftmargin=2em]
  \item 使用了哪些规则、模板、图像处理或学习方法；
  \item 输出状态如何提供给策略模块；
  \item 如何避免依赖测试阶段不可用的信息。
\end{itemize}

\subsection{误差与局限}

说明状态抽取中可能存在的误识别、泛化问题和对后续策略的影响。

\section{策略与规划}

\subsection{基线策略}

此处记录最初的 rule-based / search baseline，包括：
\begin{itemize}[leftmargin=2em]
  \item 如何解决 Task 1 和 Task 2；
  \item 如何表示子目标；
  \item 如何进行路径规划或动作选择。
\end{itemize}

\subsection{多任务扩展}

说明如何扩展到多房间、按钮、桥、怪物等复杂机制。

\subsection{最终提交版本}

明确区分：
\begin{itemize}[leftmargin=2em]
  \item 调试版策略；
  \item 最终提交版策略；
  \item 两者在输入信息上的区别。
\end{itemize}

\section{Lean 形式化与证明}

\subsection{形式化对象总览}

本项目对 \texttt{task\_1} 至 \texttt{task\_5} 全部完成了 Lean 环境形式化与可完成性证明，
对应文件位于 \texttt{student\_agent/lean/} 目录下：
\begin{itemize}[leftmargin=2em]
  \item \texttt{Task1EnvironmentFormalization.lean}：tile 级单房间钥匙门；
  \item \texttt{Task2EnvironmentFormalization.lean}：zone 级战斗与条件门；
  \item \texttt{Task3EnvironmentFormalization.lean}：room 级三房间往返取钥匙；
  \item \texttt{Task4EnvironmentFormalization.lean}：room 级旋转桥与隐藏宝箱；
  \item \texttt{Task5EnvironmentFormalization.lean}：room 级四房间按钮门与全宝箱收集。
\end{itemize}

五个文件均通过 Lean LSP 诊断（零错误），并按照统一的设计原则建模：
\begin{enumerate}[leftmargin=2em]
  \item \textbf{抽象层级与 Python 行为对齐}：根据每个任务的关键机制选择
  tile / zone / room 抽象层级，并与
  \texttt{movement.py}、\texttt{interactions.py}、\texttt{combat.py}
  及 \texttt{room\_*.json} 中的语义保持一致；
  \item \textbf{保留约束而不照搬实现}：刻意忽略像素、碰撞盒、怪物 AI、击退等低层细节，
  只保留影响规划正确性的状态字段、前置条件与事件触发；
  \item \textbf{统一证明骨架}：每个任务都给出
  “局部前置/后置引理 $\rightarrow$ 必要性引理 $\rightarrow$ \texttt{Reachable} 归纳不变式
  $\rightarrow$ \texttt{completion\_postcondition} $\rightarrow$ witness 计划 $\rightarrow$ 可完成性”
  的完整证明链。
\end{enumerate}

下面分任务展开。其中 \texttt{task\_4} 因机制最复杂（动态对象状态、条件门、
关键道具获取和事件驱动显现），作为重点展开对象。

\subsection{Task 1：单房间钥匙门形式化}

\texttt{task\_1} 是最简单的单房间取钥匙走北门任务。
Lean 文件以 tile 级建模，对应 \texttt{room\_001.json} 的 8 行 $\times$ 10 列网格。

\textbf{状态与动作}：
\begin{itemize}[leftmargin=2em]
  \item \texttt{EnvState}：\texttt{x, y : Nat}（tile 坐标）、\texttt{keys}、
  \texttt{chestOpened}、\texttt{completed}；
  \item \texttt{Action}：\texttt{up / down / left / right / openChest / goNorth}；
  \item \texttt{isWalkable} 严格复刻 \texttt{room\_001.json} 中的墙体布局
  （row 2 的 \texttt{\#\#..\#\#\#\#\#\#}、row 5 的 \texttt{\#\#\#\#\#\#\#...}）；
  \item \texttt{isAdjacentToChest} 使用 Manhattan 距离 $\leq 1$，
  对齐 \texttt{state.py:is\_adjacent} 的语义（包含距离 0）。
\end{itemize}

\textbf{关键约束}：
\begin{itemize}[leftmargin=2em]
  \item \texttt{openChest} 触发条件为邻接 chest tile 且未开过，开箱后 \texttt{keys + 1}；
  \item \texttt{goNorth} 仅在 \texttt{(4, 0)} 且 \texttt{keys > 0} 时触发，
  成功后消耗一把钥匙并将 \texttt{completed} 置真（对应 \texttt{consume\_key: true}）。
\end{itemize}

\begin{longtable}{p{0.30\textwidth}p{0.28\textwidth}p{0.32\textwidth}}
\toprule
定理名称 & 形式化对象 & 结论说明 \\
\midrule
\endhead
\texttt{step\_preserves\_walkable} & 安全性 &
任意动作后玩家仍位于可走 tile 上。 \\
\texttt{openChest\_only\_at\_chest} & 开箱前置条件 &
不邻接 chest 时 \texttt{openChest} 不改变状态。 \\
\texttt{openChest\_at\_chest\_increases\_keys} & 开箱后置条件 &
邻接且未开过时，\texttt{keys} 增加 1。 \\
\texttt{goNorth\_only\_at\_exit} & 北门前置条件 &
不在 \texttt{(4, 0)} 时 \texttt{goNorth} 无效。 \\
\texttt{goNorth\_without\_key} & 北门前置条件 &
在出口但无钥匙时 \texttt{goNorth} 无效。 \\
\texttt{goNorth\_at\_exit\_with\_key\_completes} & 北门后置条件 &
在出口且持钥匙时 \texttt{completed} 置真。 \\
\texttt{only\_goNorth\_sets\_completed} & 必要性引理 &
\texttt{goNorth} 是唯一能把 \texttt{completed} 从假翻为真的动作。 \\
\texttt{witnessPlan\_reaches\_goal} & 执行轨迹 &
22 步 witness 计划可达目标。 \\
\texttt{task1\_completable} & 可完成性 &
\texttt{task\_1} 在该模型下可完成。 \\
\bottomrule
\end{longtable}

\textbf{witness 计划}：tile 级路径
\[
(4,6) \xrightarrow{\text{right}\times 3} (7,6)
  \xrightarrow{\text{up}\times 3} (7,3)
  \xrightarrow{\text{left}\times 7} (0,3)
  \xrightarrow{\text{openChest}} \cdot
  \xrightarrow{\text{right}\times 3, \text{up}\times 3, \text{right}} (4,0)
  \xrightarrow{\text{goNorth}} \text{goal}.
\]
由于 22 步计划超出 \texttt{decide} 递归深度，按项目约定使用 \texttt{native\_decide} 验证
\texttt{witnessPlan\_executes\_to\_finalState}。

\subsection{Task 2：战斗与条件门形式化}

\texttt{task\_2} 在单房间内引入 chaser 怪物、HP、陷阱与西向条件门。
为突出战斗与条件门语义，采用 zone 级抽象。

\textbf{状态与动作}：
\begin{itemize}[leftmargin=2em]
  \item \texttt{Zone}：\texttt{safeCenter / nearMonster / nearChest / westExit / topTrap / bottomTrap}；
  \item \texttt{EnvState}：\texttt{zone}、\texttt{hp}、\texttt{monsterAlive}、\texttt{monsterHp}、
  \texttt{chestOpened}、\texttt{keys}、\texttt{completed}；
  \item \texttt{Action}：\texttt{moveTo z / attack / openChest / useExit / wait}。
\end{itemize}

\textbf{关键约束}：
\begin{itemize}[leftmargin=2em]
  \item \texttt{attack} 仅在 \texttt{nearMonster}、怪物存活且 \texttt{hp > 1} 时生效；
  \texttt{monsterHp} 为 1 时击杀（置 \texttt{monsterAlive := false}），
  否则仅 \texttt{monsterHp - 1}；
  \item \texttt{useExit} 在 \texttt{westExit} 且怪物已死且 \texttt{keys > 0} 时触发，
  成功后\textbf{不消耗}钥匙（对应 \texttt{requires} 中未设 \texttt{consume\_key}）；
  \item 走入 \texttt{topTrap / bottomTrap} 时 HP 减 1 并被重置回 \texttt{safeCenter}。
\end{itemize}

\begin{longtable}{p{0.32\textwidth}p{0.26\textwidth}p{0.32\textwidth}}
\toprule
定理名称 & 形式化对象 & 结论说明 \\
\midrule
\endhead
\texttt{trap\_move\_reduces\_hp} & 陷阱语义 &
踩陷阱后 \texttt{hp} 减少 1。 \\
\texttt{safe\_move\_preserves\_*} & 安全移动不变式 &
非陷阱 \texttt{moveTo} 保持 HP/怪物/宝箱/钥匙/完成状态不变。 \\
\texttt{attack\_reduces\_monsterHp\_when\_adjacent} & 战斗语义 &
邻接且存活且 HP 充足时，\texttt{monsterHp} 减 1。 \\
\texttt{attack\_kills\_monster\_when\_monsterHp\_one} & 战斗语义 &
\texttt{monsterHp = 1} 时一击击杀。 \\
\texttt{attack\_no\_effect\_*} & 战斗前置条件 &
不在 \texttt{nearMonster}、怪物已死或 HP 过低时攻击无效。 \\
\texttt{exit\_blocked\_if\_monster\_alive} & 条件门 &
怪物存活时 \texttt{useExit} 无效。 \\
\texttt{exit\_blocked\_without\_key} & 条件门 &
无钥匙时 \texttt{useExit} 无效。 \\
\texttt{exit\_succeeds\_after\_requirements} & 条件门后置 &
怪物已死且持钥匙时 \texttt{completed} 置真。 \\
\texttt{only\_useExit\_sets\_completed} & 必要性引理 &
\texttt{useExit} 是唯一完成动作。 \\
\texttt{completion\_postcondition} & 完成后置 &
任意可达完成态满足怪物已死 $\wedge$ 宝箱已开。 \\
\texttt{hp\_non\_increasing} & 安全性 &
任意动作后 HP 不增。 \\
\texttt{witnessPlan\_reaches\_goal} & 执行轨迹 &
7 步 witness 计划可达目标。 \\
\texttt{task2\_completable} & 可完成性 &
\texttt{task\_2} 在该模型下可完成。 \\
\bottomrule
\end{longtable}

\textbf{witness 计划}：zone 级路径
\[
\text{safeCenter} \to \text{nearMonster}
  \xrightarrow{\text{attack}\times 2} \cdot
  \to \text{nearChest} \xrightarrow{\text{openChest}} \cdot
  \to \text{westExit} \xrightarrow{\text{useExit}} \text{goal}.
\]

\subsection{Task 3：多房间往返形式化}

\texttt{task\_3} 是首个多房间任务：从 \texttt{start\_room} 出发向西穿过
\texttt{monster\_hall} 到达 \texttt{key\_room}，取钥匙后返回
\texttt{start\_room} 走东向锁定出口。采用 room 级抽象。

\textbf{状态与动作}：
\begin{itemize}[leftmargin=2em]
  \item \texttt{Room}：\texttt{startRoom / monsterHall / keyRoom}；
  \item \texttt{EnvState}：\texttt{room}、\texttt{keys}、\texttt{keyChestOpened}、
  \texttt{hallMonsterAlive}、\texttt{completed}；
  \item \texttt{Action}：\texttt{goWest / goEast / openChest / openLockedExit / wait}。
\end{itemize}

\textbf{关键约束}：
\begin{itemize}[leftmargin=2em]
  \item \texttt{goWest}：\texttt{startRoom $\to$ monsterHall $\to$ keyRoom}；
  \texttt{goEast}：反向；\texttt{keyRoom} 的 \texttt{goWest} 与
  \texttt{startRoom} 的 \texttt{goEast} 无效；
  \item \texttt{openChest} 仅在 \texttt{keyRoom} 且未开过时生效，开箱后 \texttt{keys + 1}；
  \item \texttt{openLockedExit} 仅在 \texttt{startRoom} 且 \texttt{keys > 0} 时触发，
  成功后消耗钥匙并完成（对应 \texttt{consume\_key: true}）。
\end{itemize}

\begin{longtable}{p{0.34\textwidth}p{0.24\textwidth}p{0.32\textwidth}}
\toprule
定理名称 & 形式化对象 & 结论说明 \\
\midrule
\endhead
\texttt{west\_then\_west\_reaches\_keyRoom} & 导航语义 &
从 \texttt{startRoom} 连续两次 \texttt{goWest} 可达 \texttt{keyRoom}。 \\
\texttt{openChest\_only\_in\_keyRoom} & 开箱前置 &
不在 \texttt{keyRoom} 时 \texttt{openChest} 无效。 \\
\texttt{openChest\_in\_keyRoom\_increases\_keys} & 开箱后置 &
在 \texttt{keyRoom} 且未开过时 \texttt{keys} 增 1。 \\
\texttt{openLockedExit\_blocked\_without\_key} & 锁门前置 &
在 \texttt{startRoom} 但无钥匙时 \texttt{openLockedExit} 无效。 \\
\texttt{openLockedExit\_consumes\_key\_and\_completes} & 锁门后置 &
持钥匙时消耗一把钥匙并完成。 \\
\texttt{hallMonsterAlive\_never\_changes} & 不变式 &
Task 3 无攻击动作，\texttt{hallMonsterAlive} 恒不变。 \\
\texttt{only\_openLockedExit\_sets\_completed} & 必要性引理 &
\texttt{openLockedExit} 是唯一完成动作。 \\
\texttt{keyChestOpened\_or\_keys\_zero} & 归纳不变式 &
任意可达状态：宝箱已开 $\vee$ 钥匙数为 0。 \\
\texttt{completion\_postcondition} & 完成后置 &
任意可达完成态满足 \texttt{keyChestOpened = true}。 \\
\texttt{witnessPlan\_reaches\_goal} & 执行轨迹 &
6 步 witness 计划可达目标。 \\
\texttt{task3\_completable} & 可完成性 &
\texttt{task\_3} 在该模型下可完成。 \\
\bottomrule
\end{longtable}

\textbf{witness 计划}：
\[
\text{startRoom} \xrightarrow{\text{goWest}\times 2} \text{keyRoom}
  \xrightarrow{\text{openChest}} \cdot
  \xrightarrow{\text{goEast}\times 2} \text{startRoom}
  \xrightarrow{\text{openLockedExit}} \text{goal}.
\]
由于计划简单，\texttt{witnessPlan\_executes\_to\_finalState} 用 \texttt{rfl} 直接验证。

\subsection{Task 4：旋转桥与隐藏宝箱形式化}

\texttt{task\_4} 引入旋转桥、条件门、装备获取与事件驱动显现机制，是机制最复杂的一关，
也是本报告的重点展开对象。对应 Lean 文件为
\texttt{student\_agent/lean/Task4EnvironmentFormalization.lean}，
与 Python 环境中以下模块对齐：
\texttt{movement.py}（门与房间切换）、
\texttt{interactions.py}（开关与宝箱）、
\texttt{combat.py}（击杀怪物后揭示宝箱）。

\textbf{具体抽象}：
\begin{itemize}[leftmargin=2em]
  \item \textbf{房间状态}：使用 \texttt{Room} 枚举描述五个房间
  \texttt{west / center / north / east / south}。
  \item \textbf{动态对象状态}：使用 \texttt{BridgeState} 描述中央桥的三种连通状态
  \texttt{westToNorth / westToEast / westToSouth}。
  \item \textbf{玩家资源与关键事件}：记录
  \texttt{keys}、\texttt{hasSword}、\texttt{hasShield}，
  以及北房宝箱、东房宝箱、南房守卫和中央终点宝箱的状态。
  \item \textbf{动作}：抽象为房间级动作，包括转桥、进出中心房、开宝箱、攻击和开启最终宝箱。
  \item \textbf{转移函数}：用确定性函数 \texttt{step} 描述状态更新。
  它保留了三类关键约束：
  东门需要钥匙但不消耗钥匙（对应 \texttt{consume\_key: false}）；
  桥状态决定中心房可达方向；
  南房守卫被击杀后才会揭示中央最终宝箱。
  \item \textbf{目标谓词}：\texttt{GoalReached} 定义为最终宝箱已开启。
\end{itemize}

\begin{longtable}{p{0.28\textwidth}p{0.28\textwidth}p{0.34\textwidth}}
\toprule
定理名称 & 形式化对象 & 结论说明 \\
\midrule
\endhead
\texttt{rotateBridge\_three\_cycle} & 桥状态机 &
证明西房开关连续触发三次后，桥状态回到初始方向。 \\
\texttt{goEast\_blocked\_without\_key} & 东门条件 &
证明当玩家位于中心房、桥已连到东侧但没有钥匙时，\texttt{goEast} 不产生状态变化。 \\
\texttt{goEast\_succeeds\_with\_key} & 东门条件 &
证明在中心房、桥连东侧且已持钥匙时，\texttt{goEast} 能把玩家带入东房。 \\
\texttt{goEast\_preserves\_keys} & 东门不变式 &
证明 \texttt{goEast} 不消耗钥匙（Python \texttt{consume\_key: false} 对齐）。 \\
\texttt{attack\_without\_sword\_has\_no\_effect} & 战斗语义 &
证明没有剑时，在南房执行攻击不会击杀守卫，也不会揭示最终宝箱。 \\
\texttt{attack\_with\_sword\_reveals\_final\_chest} & 战斗与事件触发 &
证明持剑击败南房守卫后，\texttt{guardianAlive} 变为假，且 \texttt{finalChestVisible} 变为真。 \\
\texttt{openFinalChest\_requires\_visibility} & 终点交互语义 &
证明若最终宝箱尚未显现，则在中心房尝试开启最终宝箱不会成功。 \\
\texttt{only\_openFinalChest\_sets\_finalChestOpened} & 必要性引理 &
\texttt{openFinalChest} 是唯一完成动作。 \\
\texttt{northChestOpened\_or\_keys\_zero} & 归纳不变式 &
任意可达状态：北房宝箱已开 $\vee$ 钥匙数为 0。 \\
\texttt{finalChestVisible\_implies\_guardianDead} & 归纳不变式 &
任意可达状态：终点宝箱显现 $\Rightarrow$ 守卫已死。 \\
\texttt{completion\_postcondition} & 完成后置 &
任意可达完成态满足 \texttt{guardianAlive = false}。 \\
\texttt{witnessPlan\_reaches\_goal} & 执行轨迹 &
给出一条具体计划，证明从初始状态出发可以完成“取钥匙 $\rightarrow$ 拿剑 $\rightarrow$ 击杀守卫 $\rightarrow$ 开最终宝箱”的任务链。 \\
\texttt{task4\_completable} & 可达性 &
由上述执行轨迹推出：在该形式化模型下，\texttt{task\_4} 是可完成的。 \\
\bottomrule
\end{longtable}

\textbf{witness 计划}：17 步房间级路径
\[
\begin{aligned}
&\text{west} \xrightarrow{\text{goCenter, goNorth, openChest}} \text{north}
  \xrightarrow{\text{goCenter, goWest, toggleBridge}} \text{west} \\
&\xrightarrow{\text{goCenter, goEast, openChest}} \text{east}
  \xrightarrow{\text{goCenter, goWest, toggleBridge}} \text{west} \\
&\xrightarrow{\text{goCenter, goSouth, attack}} \text{south}
  \xrightarrow{\text{goCenter, openFinalChest}} \text{goal}.
\end{aligned}
\]
\texttt{witnessPlan\_executes\_to\_finalState} 用 \texttt{rfl} 直接验证。

\subsection{Task 5：四房间按钮门与全宝箱收集形式化}

\subsubsection{环境建模}

Task 5 包含 4 个房间（\texttt{startRoom / southRoom / eastRoom / westRoom}），
关键机制为：
\begin{itemize}[leftmargin=2em]
  \item \textbf{按钮门}：\texttt{startRoom} 的南出口为条件门（\texttt{button\_pressed}），
        初始 \texttt{buttonPressed = false}；
  \item \textbf{钥匙门}：\texttt{startRoom} 的东出口为锁门（\texttt{consume\_key = true}），
        初始 \texttt{keys = 0}；
  \item \textbf{差异化宝箱}：4 个房间各有 1 个宝箱，效果不同——
        \texttt{startRoom}（金币，无副作用）、\texttt{southRoom}（钥匙，\texttt{keys + 1}）、
        \texttt{eastRoom}（回血，\texttt{hp + 1}）、\texttt{westRoom}（金币，无副作用）；
  \item \textbf{世界完成}：与 \texttt{task\_1--4} 不同，task\_5 无 \texttt{complete\_task} 出口，
        世界完成条件为 \texttt{all\_chests\_opened}（对齐 \texttt{progress.py:36}）。
\end{itemize}

\texttt{GoalReached} 定义为 \texttt{allChestsOpened}，即四个宝箱全部打开。
初始状态：\texttt{room = startRoom}，\texttt{buttonPressed = false}，\texttt{keys = 0}，\texttt{hp = 5}，
四个 \texttt{ChestOpened} 均为 \texttt{false}。

\subsubsection{状态空间与转移函数}

\begin{itemize}[leftmargin=2em]
  \item \texttt{pressButton}：仅在 \texttt{startRoom} 生效，\texttt{buttonPressed := true}；
  \item \texttt{goSouth}：需 \texttt{room = startRoom ∧ buttonPressed = true}，进入 \texttt{southRoom}；
  \item \texttt{goNorth}：需 \texttt{room = southRoom}，返回 \texttt{startRoom}；
  \item \texttt{goEast}：从 \texttt{startRoom}（\texttt{keys > 0}）进入 \texttt{eastRoom} 并消耗钥匙；
        从 \texttt{westRoom} 返回 \texttt{startRoom}（不消耗钥匙）；
  \item \texttt{goWest}：从 \texttt{startRoom} 进入 \texttt{westRoom}；从 \texttt{eastRoom} 返回 \texttt{startRoom}；
  \item \texttt{openChest}：按房间分支，每房间最多开一次，效果见上。
\end{itemize}

\subsubsection{核心定理}

\begin{table}[ht]
\centering
\caption{Task 5 核心定理一览}
\label{tab:task5-theorems}
\begin{tabular}{p{6.5cm}p{5.5cm}}
\toprule
\textbf{定理} & \textbf{作用} \\
\midrule
\texttt{pressButton\_only\_in\_startRoom} & 按钮仅起点房生效 \\
\texttt{goSouth\_blocked\_without\_button} & 条件门前置 \\
\texttt{goEast\_blocked\_without\_key} & 锁门前置 \\
\texttt{goEast\_from\_start\_consumes\_key} & 东出口消耗钥匙 \\
\texttt{openChest\_south\_grants\_key} & 南房宝箱给钥匙 \\
\texttt{openChest\_east\_heals} & 东房宝箱回血 \\
\texttt{only\_pressButton\_changes\_buttonPressed} & 必要性：按钮 \\
\texttt{non\_goEast\_non\_openChest\_preserves\_keys} & 必要性：钥匙 \\
\texttt{buttonPressed\_false\_implies\_not\_in\_south\_and\_chest\_closed} & 归纳不变式 \\
\texttt{southChestOpened\_implies\_buttonPressed} & 南宝箱开 $\Rightarrow$ 按钮已按 \\
\texttt{southChestClosed\_implies\_keys\_zero} & 钥匙来源不变式 \\
\texttt{completion\_postcondition} & 完成态 $\Rightarrow$ buttonPressed = true \\
\texttt{task5\_completable} & 可完成性 \\
\bottomrule
\end{tabular}
\end{table}

\subsubsection{归纳不变式}

\texttt{buttonPressed\_false\_implies\_not\_in\_south\_and\_chest\_closed}
是核心不变式：若 \texttt{buttonPressed = false}，则玩家不在 \texttt{southRoom} 且南宝箱未开。
其逆否命题 \texttt{southChestOpened\_implies\_buttonPressed} 说明南宝箱开启蕴含按钮已按。

\texttt{southChestClosed\_implies\_keys\_zero}
利用归纳法证明：南宝箱是唯一的钥匙来源，而东出口（唯一消耗钥匙的动作）需要 \texttt{keys > 0}，
因此在南宝箱未开之前不可能存在任何钥匙。

\subsubsection{Witness 计划}

10 步计划：
\[
\begin{aligned}
&\text{start}: \xrightarrow{\text{openChest}} \text{start}
  \xrightarrow{\text{pressButton}} \text{start}
  \xrightarrow{\text{goSouth}} \text{south} \\
&\xrightarrow{\text{openChest (keys+1)}} \text{south}
  \xrightarrow{\text{goNorth}} \text{start}
  \xrightarrow{\text{goEast (keys-1)}} \text{east} \\
&\xrightarrow{\text{openChest (hp+1)}} \text{east}
  \xrightarrow{\text{goWest}} \text{start}
  \xrightarrow{\text{goWest}} \text{west}
  \xrightarrow{\text{openChest}} \text{goal}.
\end{aligned}
\]

终态：\texttt{room = westRoom}，\texttt{buttonPressed = true}，\texttt{keys = 0}，\texttt{hp = 6}，
四个 \texttt{ChestOpened} 均为 \texttt{true}。
\texttt{witnessPlan\_executes\_to\_finalState} 用 \texttt{rfl} 直接验证。

\subsection{证明统一思路}

五个任务遵循同一套证明骨架，便于复用与审阅：

\begin{enumerate}[leftmargin=2em]
  \item \textbf{局部前置/后置引理}：对每个动作给出“何时触发”和“触发后字段如何变化”。
  这一步对应评分细则中的“状态 / 动作 / 转移是否定义完整、是否能证明基本合法性与安全性”。
  例如 Task 4 的 \texttt{goEast\_blocked\_without\_key} 直接对应 Python
  \texttt{movement.py} 对 \texttt{locked\_key} 出口的检查；
  \texttt{attack\_with\_sword\_reveals\_final\_chest} 对应
  \texttt{combat.py} 中“清怪后揭示宝箱”的事件语义；
  \texttt{rotateBridge\_three\_cycle} 对应 \texttt{interactions.py}
  中 switch 对 \texttt{center\_bridge} 的循环切换。
  \item \textbf{必要性引理}（\texttt{only\_X\_sets\_completed}）：
  证明只有某个特定动作能把完成标志从假翻为真。
  通过 \texttt{non\_X\_preserves\_completed} 这一“其它动作保持完成字段不变”的引理来支撑。
  \item \textbf{\texttt{Reachable} 归纳不变式}：定义归纳谓词
  \texttt{Reachable}，对初始状态和任意 \texttt{step} 闭合；
  证明诸如“未开箱则钥匙数为 0”“终点宝箱显现则守卫已死”等性质在所有可达状态上成立。
  \item \textbf{\texttt{completion\_postcondition}}：在 \texttt{Reachable} 归纳之上，
  证明“任意可达完成态必然满足某组前置条件”（如 Task 2 中“怪物已死 $\wedge$ 宝箱已开”、
  Task 4 中“守卫已死”）。证明中显式调用必要性引理与归纳不变式。
  \item \textbf{witness 计划}：在局部规则已经明确的基础上，构造一条显式的见证计划
  \texttt{witnessPlan}。由于 \texttt{step} 被定义为确定性函数，
  因此整条轨迹的正确性可以通过归约直接验证
  （简单计划用 \texttt{rfl}，长计划用 \texttt{native\_decide}）。
  \item \textbf{可完成性}：由 \texttt{witnessPlan\_reaches\_goal} 直接推出
  \texttt{taskN\_completable}，即“存在计划使目标谓词成立”。
\end{enumerate}

这使得“任务可完成”从口头分析转化为可检查的形式化结论，
并且完成态的必要前置条件也被显式记录下来。

\subsection{证明局限与边界}

当前环境形式化刻画的是 \texttt{task\_1} 至 \texttt{task\_5} 的
\textbf{tile / zone / room 级符号语义}，
而不是 Python 引擎的逐像素精确副本，因此存在以下边界：
\begin{itemize}[leftmargin=2em]
  \item 未形式化像素级移动、碰撞盒、路径细节与房间内精确坐标
  （\texttt{task\_1} 保留了 tile 级网格，但 \texttt{task\_2} 至 \texttt{task\_5}
  均采用更高的 zone / room 抽象）。
  \item 未形式化怪物 AI、击退、受伤时序和 shield 格挡的连续动力学。
  \item 未证明视觉模块一定能从 \texttt{obs} 中稳定恢复这些符号状态；
  该部分属于感知层，而不是当前 Lean 环境模型的覆盖范围。
  \item 当前证明的是一条显式计划的可执行性和若干局部环境性质，
  还没有证明“任意合法 planner 都能找到该计划”这类更强的完备性命题。
\end{itemize}

这一边界是有意为之：课程要求强调 Lean 证明应覆盖“可验证层”。
对于本项目而言，最具价值且最容易和实际代码保持一致的对象，
正是钥匙/门条件、战斗与事件触发、桥状态机与任务完成谓词这类符号层环境语义。

\section{实验设计与结果分析}

\subsection{实验设置}

此处记录：
\begin{itemize}[leftmargin=2em]
  \item 测试任务；
  \item 运行次数；
  \item 成功率统计方式；
  \item 使用的模型、参数或规则版本。
\end{itemize}

\subsection{实验结果}

建议用表格和图片展示：

\begin{table}[H]
\centering
\begin{tabular}{llll}
\toprule
任务 & 方法版本 & 结果指标 & 备注 \\
\midrule
Task 1 & 待补充 & 待补充 & 待补充 \\
Task 2 & 待补充 & 待补充 & 待补充 \\
Task 3 & 待补充 & 待补充 & 待补充 \\
Task 4 & 待补充 & 待补充 & 待补充 \\
Task 5 & 待补充 & 待补充 & 待补充 \\
\bottomrule
\end{tabular}
\caption{各任务实验结果汇总}
\end{table}

\subsection{结果分析}

重点分析：
\begin{itemize}[leftmargin=2em]
  \item 哪些任务成功率较高；
  \item 哪些失败来自视觉误差，哪些来自策略局限；
  \item 哪些设计有较好的泛化性；
  \item reward 和中间目标是否有帮助。
\end{itemize}

\section{项目分工与推进过程}

\subsection{成员分工}

建议写成表格：

\begin{table}[H]
\centering
\begin{tabular}{lll}
\toprule
成员 & 负责内容 & 主要产出 \\
\midrule
组长 & 待补充 & 待补充 \\
成员 1 & 待补充 & 待补充 \\
成员 2 & 待补充 & 待补充 \\
成员 3 & 待补充 & 待补充 \\
成员 4 & 待补充 & 待补充 \\
\bottomrule
\end{tabular}
\caption{小组分工}
\end{table}

\subsection{时间安排}

此处根据进度文档补充每日推进情况、关键决策与风险应对。

\section{总结与展望}

总结本文完成的工作、取得的结果、已有不足，以及如果时间更充足可以继续推进的方向。

\newpage
\bibliographystyle{plain}
\bibliography{references}

\end{document}
